\documentclass[preview]{standalone}

\usepackage{amsmath}
\usepackage{amssymb}
\usepackage{stellar}
\usepackage{definitions}

\begin{document}

\id{psicologia-modelli-comunicazione}
\genpage

\section{Modelli della comunicazione}

\begin{snippet}{comunicazione-introduzione}
    La comunicazione accompagna l'individuo per tutta l'esistenza. È un'attività
    di interazione, una necessità umana e compare fin dai primissimi giorni di
    vita attraverso segnali come pianto, sorrisi e vocalizzazioni.
\end{snippet}

\begin{snippetexample}{pianto-bambino-example}{Pianto del bambino}
    Il pianto di un bambino può avere diverse finalità comunicative: fame, noia
    o dolore specifico, come una colica. La sensibilità materna consiste nella
    capacità di sintonizzarsi con il bambino e interpretare correttamente questi
    segnali, anche quando appaiono acusticamente simili.
\end{snippetexample}

\begin{snippetdefinition}{informazione-shannon-definition}{Informazione}
    Per Shannon, l'\emph{informazione} è una grandezza fisica, osservabile e
    misurabile: una relazione tra due o più dati capace di generare ulteriori
    conoscenze.
\end{snippetdefinition}

\begin{snippetdefinition}{entropia-comunicazione-definition}{Entropia}
    L'\emph{entropia} è l'assenza di informazione, che si verifica quando il
    segnale e le interferenze alla fonte sono equiprobabili ed equiparabili.
\end{snippetdefinition}

\begin{snippet}{modello-shannon-weaver}
    Il modello di Shannon e Weaver descrive la comunicazione come trasmissione
    di un segnale da una fonte a un destinatario. Il suo requisito fondamentale
    è che fonte e destinatario condividano lo stesso codice, altrimenti
    l'informazione non può essere veicolata correttamente.
\end{snippet}

\includesnpt[width=68\%,src=/snippet/static-psychology/shannon-weaver-communication-model.jpg]{centered-img}

\begin{snippetdefinition}{codice-comunicazione-definition}{Codice}
    Il \emph{codice} è l'insieme di regole che associa in modo coerente e
    univoco gli elementi di un sistema con quelli di un altro sistema, come
    accade nella struttura sintattica o nelle associazioni semantiche.
\end{snippetdefinition}

\begin{snippetdefinition}{processo-significazione-definition}{Processo di significazione}
    Il \emph{processo di significazione} è il processo attraverso cui un segno
    rimanda a un significato e a un oggetto o evento del mondo.
\end{snippetdefinition}

\begin{snippet}{triangolo-significazione}
    Il processo di significazione comprende referente, referenza e simbolo. Il
    \emph{referente} è ciò di cui si parla, la \emph{referenza} è la
    rappresentazione mentale corrispondente e il \emph{simbolo} è l'immagine
    sonora o iconica dell'oggetto o evento.
\end{snippet}

\includesnpt[width=55\%,src=/snippet/static-psychology/triangle-of-meaning.jpg]{centered-img}

\begin{snippetdefinition}{semiotica-definition}{Semiotica}
    La \emph{semiotica} è la disciplina che studia il processo di significazione.
\end{snippetdefinition}

\begin{snippet}{peirce-segni}
    Per Peirce, il segno rimanda a qualcosa di diverso da sé. I segni possono
    essere \emph{icone}, quando si fondano su somiglianza con il referente;
    \emph{indici}, quando dipendono da connessione o analogia con l'oggetto;
    \emph{simboli}, quando la connessione con il referente è convenzionale e
    arbitraria.
\end{snippet}

\begin{snippetdefinition}{pragmatica-definition}{Pragmatica}
    La \emph{pragmatica} studia l'uso concreto di segni, parole, frasi e gesti
    in un dato contesto comunicativo.
\end{snippetdefinition}

\begin{snippetdefinition}{atto-linguistico-definition}{Atto linguistico}
    Un \emph{atto linguistico} è un'azione compiuta attraverso il linguaggio.
\end{snippetdefinition}

\begin{snippet}{atti-linguistici}
    Secondo Austin, dire qualcosa significa anche fare qualcosa. Ogni atto
    linguistico comprende un \emph{atto locutorio}, cioè il parlare stesso, un
    \emph{atto illocutorio}, cioè l'intenzione comunicativa, e un
    \emph{atto perlocutorio}, cioè gli effetti prodotti sull'interlocutore.
\end{snippet}

\begin{snippetdefinition}{principio-cooperazione-definition}{Principio di cooperazione}
    Il \emph{principio di cooperazione} afferma che l'interlocutore dovrebbe
    fornire il proprio contributo alla conversazione in modo opportuno.
\end{snippetdefinition}

\begin{snippet}{massime-grice}
    Il principio di cooperazione di Grice si articola in quattro massime:
    quantità, qualità, relazione e modo. Esse richiedono di fornire informazioni
    sufficienti, dire ciò che si ritiene vero, essere pertinenti ed esprimersi in
    modo chiaro, breve e ordinato.
\end{snippet}

\begin{snippetdefinition}{implicatura-conversazionale-definition}{Implicatura conversazionale}
    L'\emph{implicatura conversazionale} è un processo di inferenza mentale che
    permette di andare oltre il significato letterale di un enunciato e di
    comprendere l'intenzione comunicativa di chi parla.
\end{snippetdefinition}

\begin{snippetdefinition}{giochi-senza-fine-definition}{Giochi senza fine}
    I \emph{giochi senza fine} sono dinamiche comunicative e psicologiche che
    tendono a rinnovarsi continuamente perché sono alimentate dagli stessi
    interlocutori coinvolti.
\end{snippetdefinition}

\begin{snippet}{comunicazione-circolare}
    Dal punto di vista psicologico, la comunicazione è una sequenza continua e
    circolare di messaggi. Ogni scambio comunicativo è anche una risposta a un
    messaggio precedente; i conflitti possono nascere quando si segmenta in modo
    arbitrario questo processo continuo.
\end{snippet}

\includesnpt[width=52\%,src=/snippet/static-psychology/circular-communication-process.jpg]{centered-img}

\end{document}
